%\chapter{Transverse Vibrations of Strings}
%\label{transverse-vibrations-of-strings}
118. Among vibrating bodies there are none that occupy a more prominent
position than Stretched Strings. From the earliest times they have been
employed for musical purposes, and in the present day they still form the
essential parts of such important instruments as the pianoforte and the violin.
To the mathematician they must always possess a peculiar interest as the
battle-field on which were fought out the controversies of D'Alembert, Euler,
Bernoulli and Lagrange, relating to the nature of the solutions of partial
differential equations. To the student of Acoustics they are doubly important.
In consequence of the comparative simplicity of their theory, they are the
ground on which difficult or doubtful questions, such as those relating to the
nature of simple tones, can be most advantageously faced; while in the form of
a Monochord or Sonometer, they afford the most generally available means for
the comparison of pitch.

The `string' of Acoustics is a perfectly uniform and flexible filament of solid
matter stretched between two fixed points - in fact an ideal body, never
actually realized in practice, though closely approximated to by most of the
strings employed in music. We shall afterwards see how to take account of any
small deviations from complete flexibility and uniformity.

The vibrations of a string may be divided into two distinct classes, which are
practically independent of one another, if the amplitudes do not exceed certain
limits. In the first class the displacements and motions of the particles are
longitudinal, so that the string always retains its straightness. The potential
energy of a displacement depends, not on the whole tension, but on the changes
of tension which occur in the various parts of the string, due to the increased
or diminished extension. In order to calculate it we must know the relation
between the extension of a string and the stretching force. The approximate law
(given by Hooke) may be expressed by saying that the extension varies as the
tension, so that if J and / denote the natural and the stretched lengths of a
string, and 7' the tension,

ial ES (1),

where is a constant, depending on the material and the section, which may be
interpreted to mean the tension that would be necessary to stretch the string
to twice its natural length, if the law applied to so great extensions, which,
in general, it is far from doing.

119. The vibrations of the second kind are transverse; that is to say, the
particles of the string move sensibly in planes perpendicular to the line of
the string. In this case the potential energy of a displacement depends upon
the general tension, and the small variations of tension accompanying the
additional stretching due to the displacement may be left out of account. It is
here assumed that the stretching due to the motion may be neglected in
comparison with that to which the string is already subject in its position of
equilibrium. Once assured of the fulfillment of this condition, we do not, in
the investigation of transverse vibrations, require to know anything further of
the law of extension.

The most general vibration of the transverse, or lateral, kind may be resolved,
as we shall presently prove, into two sets of component normal vibrations,
executed in perpendicular planes. Since it is only in the initial circumstances
that there can be any distinction, pertinent to the question, between one plane
and another, it is sufficient for most purposes to regard the motion as
entirely confined to a single plane passing through the line of the string.

In treating of the theory of strings it is usual to commence with two
particular solutions of the partial differential equation, representing the
transmission of waves in the positive and negative directions, and to combine
these in such a manner as to suit the case of a finite string, whose ends are
maintained at rest; neither of the solutions taken by itself being consistent
with the existence of nodes, or places of permanent rest. This aspect of the
question is very important, and we shall fully consider it; but it seems
scarcely desirable to found the solution in the first instance on a property so
peculiar to a uniform string as the undisturbed transmission of waves. We will
proceed by the more general method of assuming (in conformity with what was
proved in the last chapter) that the motion may be resolved into normal
components of the harmonic type, and determining their periods and character by
the special conditions of the system.

Towards carrying out this design the first step would naturally be the
investigation of the partial differential equation, to which the motion of a
continuous string is subject. But in order to throw light on a point, which it
is most important to understand clearly, the connection between finite and
infinite freedom, and the passage corresponding thereto between arbitrary
constants and arbitrary functions, we will commence by following a somewhat
different course.

120. In Chapter 111., it was pointed out that the fundamental vibration of a
string would not be entirely altered in character, if the mass were
concentrated at the middle point. Following out this idea, we see that if the
whole string were divided into a number of small parts and the mass of each
concentrated at its centre, we might by sufficiently multiplying the number of
parts arrive at a system, still of finite freedom, but capable of representing
the continuous string with any desired accuracy, so far at least as the lower
component vibrations are concerned. If the analytical solution for any number
of divisions can be obtained, its limit will give the result corresponding to a
uniform string. This is the method followed by Lagrange.

Let 1 be the length, pl the whole mass of the string, so that p denotes the
mass per unit length, 7, the tension.

Fig. 19.

The length of the string is divided into m+1 equal parts (a), so that (1).

At the m points of division equal masses (4) are supposed concentrated, which
are the representatives of the mass of the portions (a) of the string, which
they severally bisect. The mass of each terminal portion of length 4a is
supposed to be concentrated at the final points. On this understanding, we have

(10 41) WS (2).

We proceed to investigate the vibrations' of a string, itself devoid of
inertia, but loaded at each of m points equidistant (a) from themselves and
from the ends, with a mass p.

If Wry Wy cov ui, denote the lateral displacements of the loaded points,
including the initial and final points, we have the following expressions for T
and V,

Y ae Yn eo (3) y HE Ht Wa ton + era - Ped 8,

with the conditions that y, and y, vanish. These give by Lagrange's Method the
m equations of motion,

By,+Ap, +By, =0 BY, + AY, + By, =0

By, + Av, +Bbs =0F i. (5), Bent AV nit bass = 9 where A +h, Ba Bs (6).

Supposing now that the vibration under consideration is one of normal type, we
assume that ,, y,, etc. are all proportional to

\(\cos(nt -0)\), where n remains to be determined. A and B may then be regarded
as constants, with a substitution of -n for D .

If for the sake of brevity we put the determinantal equation, which gives the
values of n , assumes the form


C, 1, 0, 0, 0 m rows 1, C1, 0, 0 0, 1, G1, Ou... 0, 0, 1, G, la... (8) 0, 0, 0,
1, Cu...

From this equation the values of the roots might be found. It may be proved
that, if \(C= 2 \cos(0)\), the determinant is equivalent to
\(\sin(m+1) + \sin(6)\); but we shall attain our object with greater ease directly from (5)
by acting on a hint derived from the known results relating to a continuous
string, and assuming for trial a particular type of vibration. Thus let a solution be
\[v.=P \sin(r-1) 8 \cos(nt-e) (9),\]
a form which secures that y,=0. In order that, may vanish, (MAL) PHS Tii (10),

where 8 is an integer. Substituting the assumed values of in the equations (5),
we find that they are satisfied, provided that
\[2B \cos(B + A =0 1  (11);\]

so that the value of n in terms of is

\[v= P, \sin(r 7 008 (1,t,) . (13),\]

where

\[r1,= 2 ae \sin(3 (m+ 1) (14),\]

and P,, c, denote arbitrary constants independent of the general constitution
of the system. The m admissible values of.n are found from (14) by ascribing to
c in succession the values 1, 2, 3...m, and are all different. If we take
s=m+1, w, vanishes, so that this does not correspond to a possible vibration.
Greater values of s give only the same periods over again. If m+1 be even, one
of the values of n-that, namely, corresponding to 9-2

3=4(m + 1),-is the same as would be found in the case of only a single load
(m=1). The interpretation is obvious. In the kind - of vibration considered
every alternate particle remains at rest; so that the intermediate ones really
move as though they were attached to the centres of strings of length 2a,
fastened at the ends.

The most general solution is found by putting together all the possible
particular solutions of normal type

\[Vaz f \sin(nt -) (15),\]

and, by ascribing suitable values to the arbitrary constants, can be identified
with the vibration resulting from arbitrary initial circumstances,

. Let x denote the distance of the particle r from the end of the string, so
that (r-1)a=a; then by substituting for e and a from (1) and (2), our solution
may be written,

\[(c) =P, \sin(s) - \cos(b-6),\]

(16), In order to pass to the case of a continuous string, we have only to put
m infinite. The first equation retains its form, and specifies the displacement
at any point a. The limiting form of the second is simply

sr /T a= T p (18), whence for the periodic time, Q2r 2 /p T= yn = rs T.  (19)

The periods of the component tones are thus aliquot parts of that of the
gravest of the series, found by putting s=1. The whole motion is in all cases
periodic; and the period is 27 f.

This statement, however, must not be understood as excluding a shorter period;
for in particular cases any number of the lower components may be absent. All
that is asserted is that the above-mentioned interval of time is sufficient to
bring about a complete recurrence. We defer for the present any further
discussion of the important formula (19), but it is interesting to observe the
approach to a limit in (17), as m is made successively greater and greater. For
this purpose it will be sufficient to take the gravest tone for which s=1, and
accordingly to trace the variation of

\[2(m+1) \sin(Ea 2 (m+ 1) ,\]

The following are a series of simultaneous values of the function and variable:

It will be seen that for very moderate values of m the limit is closely
approached. Since m is the number of (moveable) loads, the case m=1 corresponds
to the problem investigated in Chapter II., but in comparing the results we
must remember that we there supposed the whole mass of the string to be
concentrated at the centre. In the present case the load at the centre is only
half as great; the remainder being supposed concentrated at the ends, where it
is without effect.

From the fact that our solution is general, it follows that any initial form of
the string can be represented by

\[yp (e)=  (P \cos(e), \sin(s) TF\]

(20). And, since any form possible for the string at all may be regarded as
initial, we infer that any finite single valued function of x, which vanishes
at = 0 and =J, can be expanded within those limits in a series of sines of se
and its multiples-which is a case of Fourier's theorem. We shall presently shew
how the more general form can be deduced.

121. We might now determine the constants for a continuous string by
integration as in Article 93, but it is instructive to solve the problem first
in the general case (m finite), and afterwards to proceed to the limit.. The
initial conditions are

\[( )=A \sin(4)+ A \sin(27) F +...4 A \sin(m), (2a) =A \sin(27) + A \sin(47)\]

See

\[+... +A, \sin(2m) 7,\]

\[(ma) =A, \sin(m)= + A, \sin(2m)  +... +4, \sin(mm);\]
where, for brevity,
A,=P,

\[\cos(c),, and (a), (2a) (ma)\]
are the initial displacements of the m
particles.

To determine any constant Ae multiply the first equation by sins 7 the second
by

\[\sin(2s) 7\]

, etc., and add the results. Then, by Trigonometry, the coefficients of all the
constants, except A,, vanish, while that of 4,=4(m+1)\footnote{Todhunter's
Integral Calculus, p. 267}. Hence

\[As=(2/(m+1)) sum \psi (ra) \sin(rs  \pi a /l (1).\]

We need not stay here to write down the values of B, (equal to P,

\[\sin(e)\]
as depending on the initial velocities When a becomes infinitely
small, ra under the sign of summation ranges by infinitesimal steps from zero
to 7. At the same time 1/(m+1) = a/l, so that writing ra = x, a = dx, we have
ultimately

\[As=(2/(m+1)) sum \psi (ra) \sin(rs  \pi a /l (2).\]

expressing A, in terms of the initial displacements.

122. We will now investigate independently the partial differential equation
governing the transverse motion of a perfectly flexible string, on the
suppositions (1) that the magnitude of the tension may be considered constant,
(2) that the square of the inclination of any part of the string to its initial
direction may be neglected. As before, p denotes the linear density at any
point, and 7', is the constant tension. Let rectangular coordinates be taken
parallel, and perpendicular to the string, so that gives the equilibrium and 2,
y,z the displaced position of any particle at time c. The forces acting on the
element dz are the tensions at its two ends, and any impressed forces

\[Y \rho dx, Z \rho dx.\]

By D'Alembert's Principle these form an equilibrating system with the reactions
against acceleration,

p i -p a

At the point x the components of tension are

if the squares of =, ds be neglected; so that the forces acting on the element
dx arising out of the tension are

Hence for the equations of motion,

from which it appears that the dependent variables y and z are altogether
independent of one another.

The student should compare these equations with the corresponding equations of
finite differences in Article 120. The latter may be written

Now in the limit, when a becomes infinitely small,

( - a) +(e + a) - B(x) = Y'"() a, while
\[\mu = \rho a;\]
and the equation
assumes ultimately the form agreeing with (1).

In like manner the limiting forms of (3) and (4) of Article 120 are

T=1p (GB) ae (2),

V=47, f (q)'ae (3), which may also be proved directly.

The first is obvious from the definition of T. To prove the second, it is
sufficient to notice that the potential energy in any configuration is-the work
required to produce the necessary stretching against the tension 7,. Reckoning
from the configuration of equilibrium, we have

v=7 G1 ae; and, so far as the third power of dy

123. In most of the applications that we shall have to make, the density p is
constant, there are no impressed forces, and the motion may be supposed to take
place in one plane. We may then conveniently write

and the differential equation is expressed by

If we now assume that y varies as \(\cos(mat),\) our equation becomes

(3) of which the most general solution is

\[y = (A \sin(max) + C \cos(mx) \cos(MAb (4).\]

This, however, is not the most general harmonic motion of the period in
question. In order to obtain the latter, we must assume;

\[Y = Y \cos(MAL) + Y, \sin(MAL),\]

(5), where y,, , are functions of a, not necessarily the same. On substitution
in (2) it appears that y, and y, are subject to equations of the form (3), so
that finally

\[y =(A \sin(mx) + C \cos(mz) \cos(mat) + (B a \sin(ma) + D \cos(mz) \sin(mat);\]

(6) an expression containing four arbitrary constants, For any continuous
length of string satisfying without interruption the differential equation,
this is the most general solution possible, under the condition that the motion
at every point shall be simple harmonic.  But whenever the string forms part of
a system vibrating freely and without dissipation, we know from former chapters
that all parts are simultaneously in the same phase, which requires that

A:B = C:D (7);

and then the most general vibration of simple harmonic type is

\[y = ( \alpha \sin(mx) + \beta \cos(mx) \cos(mat - \epsilon ) (8).\]

124. The most simple as well as the most important problem connected with our
present subject is the investigation of the free vibrations of a finite string
of length / held fast at both its ends. If we take the origin of at one end,
the terminal conditions are that when 2=0, and when =1, y vanishes for all
values of c. The first requires that in (6) of Article 123

C=0, D=0 (1);

and the second that

sin(ml) = 0 (2),

or that \(ml=s \pi\) , where s is an integer. We learn that the only harmonic
vibrations possible are such as make

\[m = s \pi /l (3),\]

and then

\[y = \sin(s \pi x/l)  (A \cos(s \pi a t /l) +  B \sin(s \pi a t/l)) (4).\]

Now we know a priori that whatever the motion may be, it can be represented as
a sum of simple harmonic vibrations, and we therefore conclude that the most
general solution for a string, fixed at 0 and l, is

\[y= sum \sin(Tr) (4 \cos() B \sin().\]

(5). The slowest vibration is that corresponding to s=1. Its period \(\tau_1\)
is given by

\[\tau_1 =  2 l/a = 2 l \sqrt{ \rho / T 1 }\]

(6). The other components have periods which are aliquot parts of \(\tau_1\)

\[\tau s = \tau_1 / s  (7);\]
so that, as has been already stated, the whole
motion is under all circumstances periodic in the time \(\tau_1.\) The sound
emitted constitutes in general a musical note, according to our definition of
that term, whose pitch is fixed by \(\tau_1\), the period of its gravest
component. It may happen, however, in special cases that the gravest vibration
is absent, and yet that the whole motion is not periodic in any shorter time.
This condition of things occurs, if A, +B, vanish, while, for example, 4,'+B,
and A, +B, are finite. In such cases the sound could hardly be called a note;
but it usually happens in practice that, when the gravest tone is absent, some
other takes its place in the character of fundamental, and the sound still
constitutes a note in the ordinary sense, though, of course, of elevated pitch.
A simple case is when all the odd components beginning with the first are
missing. The whole motion is then periodic in the time 7,, and if the second
component be present, the sound presents nothing unusual.

The pitch of the note yielded by a string (6), and the character of the
fundamental vibration, were first investigated on mechanical principles by
Brook Taylor in 1715; but it is to Daniel Bernoulli (1755) that we owe the
general solution contained in (5). He obtained it, as we have done, by the
synthesis of particular solutions, permissible in accordance with his Principle
of the Coexistence of Small Motions. In his time the generality of the result
so arrived at was open to question; in fact, it was the opinion of Euler, and
also, strangely enough, of Lagrange', that the series of sines in (5) was not
capable of representing an arbitrary function; and Bernoulli's argument on the
other side, drawn from the infinite number of the disposable constants, was
certainly inadequate.

Dangling\footnote{ See Riemann's Partielle Differential Gleichungen, Article 78.}

2 Dr Young, in his memoir of 1800, seems to have understood this matter quite
correctly. He says, `\,` At the same time, as M. Bernoulli has justly observed,
since every figure may be infinitely approximated, by considering its ordinates
as composed of the ordinates of an infinite number of trochoids of different
magnitudes, it may be demonstrated that all these constituent curves would
revert to their initial state, in the same time that a similar chord bent into
a trochoidal curve would perform a single vibration; and this is in some
respects a convenient and compendious method of considering the problem.''

124. . MERSENNE'S LAWS. 139

Most of the laws embodied in Taylor's formula (6) had been discovered
experimentally long before (1636) by Mersenne. They may be stated thus:

(1) For a given string and a given tension, the time varies as the length.

This is the fundamental principle of the monochord, and appears to have been
understood by the ancients'.  Aristotle
\begin{quote} knew that a pipe or a
chord of double length produced a sound of which the vibrations occupied a
double time; and that the properties of concords depended on the proportions of
the times occupied by the vibrations of the separate sounds.
\end{quote}
Young's Lectures on Natural Philosophy, Vol. 1., p. 404.

(2) When the length of the string is given, the time varies inversely as the
square root of the tension.

(3) Strings of the same length and tension vibrate in times, which are
proportional to the square roots of the linear density. These important results
may all be obtained by the method of dimensions, if it be assumed that +
depends only on J, p, and Z".

For, if the units of length, time and mass be denoted respectively by Z , 7; ,
2 , the dimensions of these symbols are given by

b= L , p= ML" . T,= MLT, and thus (see Article 52) the only combination of them
capable of representing a time is 7, pt J. The only thing left undetermined is
the numerical factor.

125. Mersenne's laws are exemplified in all stringed instruments. In playing
the violin different notes are obtained from the same string by shortening its
efficient length. In tuning the violin or the pianoforte, an adjustment of
pitch is effected with a constant length by varying the tension; but it must be
remembered that p is not quite invariable.

To secure.a prescribed pitch with a string of given material, it is requisite
that one relation only be satisfied between the length, the thickness, and the
tension; but in practice there is usually no great latitude. The length is
often limited by considerations of convenience, and its curtailment cannot
always be compensated by an increase of thickness, because, if the tension be
not increased proportionally to the section, there is a loss of flexibility,
while if the tension be so increased, nothing is effected towards lowering the
pitch. The difficulty is avoided in the lower strings of the pianoforte and
violin by the addition of a coil of fine wire, whose effect is to impart
inertia without too much impairing flexibility.

For quantitative investigations into the laws of strings, the sonometer is
employed. By means of a weight hanging over a pulley, a catgut, or a metallic
wire, is stretched across two bridges mounted on a resonance case. A moveable
bridge, whose Position is estimated by a scale running parallel to the wire,
gives the means of shortening the efficient portion of the wire to any desired
extent. The vibrations may be excited by plucking, as in the harp, or with a
bow (well supplied with rosin), as in the violin.

If the moveable bridge be placed half-way between the fixed ones, the note is
raised an octave; when the string is reduced to one-third, the note obtained is
the twelfth.

By means of the law of lengths, Mersenne determined for the first time the
frequencies of known musical notes. He adjusted the length of a string until
its note was one of assured position in the musical scale, and then prolonged
it under the same tension until the vibrations were slow enough to be counted.

For experimental purposes it is convenient to have two, or more, strings
mounted side by side, and to vary in turn their lengths, their masses, and the
tensions to which they are subjected. Thus in order that two strings of equal
length may yield the interval of the octave, their tensions must be in the
ratio of 1: 4, if the masses be the same; or, if the tensions be the same, the
masses must be in the reciprocal ratio.

The sonometer is very useful for the numerical determination of pitch.  By
varying the tension, the string is tuned to unison with a fork, or other
standard of known frequency, and then by adjustment of the moveable bridge, the
length of the string is determined, which vibrates in unison with any note
proposed for measurement. The law of lengths then gives the means of effecting
the desired comparison of frequencies.

Another application by Scheibler to the determination of absolute pitch is
important. The principle is the same as that explained in Chapter 1, and the
method depends on deducing the absolute pitch of two notes from a knowledge of
both the ratio and the difference of their frequencies.  The lengths of the
sonometer string when in unison with a fork, and when giving with it four beats
per second, are carefully measured. The ratio of the lengths is the inverse
ratio of the frequencies, and the difference of the frequencies is four. From
these data the absolute pitch of the fork can be calculated.


The pitch of a string may be calculated also by Taylor's formula from the
mechanical elements of the system, but great precautions are necessary to
secure accuracy. The tension is produced by a weight, whose mass (expressed
with the same unit as p) may be called P; so that 7, = gP, where g = 82:2, if
the units of length and time be the foot and the second. In order to secure
that the whole tension acts on the vibrating segment, no bridge must be
interposed, a condition only to be satisfied by suspending the string
vertically. After the weight is attached, a portion of the string is isolated
by clamping it firmly at two points, and the length is measured. The mass of
the unit of length p refers to the stretched state of the string, and may be
found indirectly by observing the elongation due to a tension of the same order
of magnitude as 7, and calculating what would be produced by 7, according to
Hooke's law, and by weighing a known length of the string in its normal state.
After the clamps have been secured great care is required to avoid fluctuations
of temperature, which would seriously influence the tension. In this way
Seebeck obtained very accurate results.

126. When a string vibrates in its gravest normal mode, the excursion is at any
moment proportional to \(\sin()=\), increasing numerically from either end
towards the centre; no intermediate point of the string remains permanently at
rest. But it is otherwise in the case of the higher normal components. Thus, if
the vibration be of the mode expressed b PY

\[y= \sin(o) + A cos\]
the excursion is proportional to sin, which vanishes at
s-1 points, dividing the string into s equal parts. These points of no motion
are called nodes, and may evidently be touched or held fast without in any way
disturbing the vibration. - The production of `harmonics' by lightly touching
the string at the points of aliquot division is a well-known resource of the
violinist. All component modes are excluded which have not a node at the point
touched; so that, as regards pitch, the effect is the same as if the string
were securely fastened there.

127. The constants, which occur in the general value of y, Article 124, depend
on the special circumstances of the vibration, and may be expressed in terms of
the initial values of y and .

Putting c = 0, we find

\[=8tp Ay \sin(2), j= TE R7 oBy \sin(7)... 0)\]

Multiplying by \(\sin(r)\) and integrating from 0 to J, we obtain

\[=F) ssi 2 ae; B= =  jp \sin(7) de\]

(2). These results exemplify Stokes' law, Article 95; for that part of y, which
depends on the initial velocities, is

and from this the part depending on initial displacements may be inferred, by
differentiating with respect to the time, and substituting y, for

When the condition of the string at some one moment is thoroughly known, these
formulae allow us to calculate the motion for all subsequent time. For example,
let the string be initially at rest, and so displaced that it forms two sides
of a triangle. Then B=0; and

Fig. 20.

on integration. We see that .A, vanishes, if \(\sin(0),\) that is, if there be
a node of the component in question situated at P. A more comprehensive view of
the subject will be afforded by another mode of solution to be given presently.

% Many articles refer back to this one.
128. In the expression for y the coefficients of \(\sin()\) are the normal
coordinates of Chapters IV., and V. We will denote them therefore by c,, 80
that the configuration and motion of the system at any instant are defined by
the values of c, and c, according to the equations

\[y= \sin()+g \sin(7) + \sin()+ (1).\]
\[y= c \sin(Te) 4 \sin()= Pb the \sin()+ wes\]

We proceed to form the expressions for J' and V, and thence to deduce the
normal equations of vibration.

For the kinetic energy,

Tip y dx= t of, 3 at dx ein dx

\[= t of  nt oy \sin(in) sd 9\]
the product of every pair of terms varying by
the general property of normal coordinates. Hence

\[T= (1/2) \rho \]
(2).  In like manner,
\[V=\frac{1}{2} T 1.\]
(3).  These
expressions do not presuppose any particular motion, either natural, or
otherwise; but we may apply them to calculate the whole energy of a string
vibrating naturally, as follows: -If M be the whole mass of the string ( \(\rho
l\) ), and its equivalent (a2 p) be substituted for T1, we find for the sum of
the energies,

\[T+V= (1/4) M (4)\]
or, in terms of As, and Bs, of Article 126,

\[T+V=  M (8)\]

If the motion be not confined to the plane of ay, we have merely to add the
energy of the vibrations in the perpendicular plane.

Lagrange's method gives mae the equation of motion \(\phi +\phi (6),\)

which has been already considered in Article 66. If c, and c, be the initial
values of c and c, the general solution is

\[=b- \sin(nt + c \cos(nt) +2 \sin(n E) (7),\]
where n is written for T

By definition , is such that. , 8c, represents the work done by the impressed
forces on the displacement 5c,. Hence, if the force acting at time c on an
element of the string p dz be p Ydz,

In these equations c, is a linear quantity, as we see from (1); and , is
therefore a force of the ordinary kind.

129. In the applications that we have to make, the only impressed force will be
supposed to act in the immediate neighbourhood of one point 2=b, and may
usually be reckoned as a whole, so that

Dangling\footnote{ The observation that a harmonic is not generated, when one of its nodal
points is plucked, is due to Young.}

If the point of application of the force coincide with a node of the mode (s),
,=0, and we learn that the force is altogether without influence on the
component in question. This principle is of great importance; it shews, for
example, that if a string be at rest in its position of equilibrium, no force
applied at its centre, whether in the form of plucking, striking, or bowing,
can generate any of the even normal components'. If after the operation of the
force, its point of application be damped, as by touching it with the finger,
all motion must forthwith cease; for those components which have not a node at
the point in question are stopped by the damping, and those which have, are
absent from, the beginning', More generally, by damping any point of a sounding
string, we stop all the component vibrations which have not, and leave entirely
unaffected those which have a node at the point touched.


The case of a string pulled aside at one point and afterwards let go from rest
may be regarded as included in the preceding statements. The complete solution
may be obtained thus. Let the motion commence at the time c=0; from which
moment ,=0. The value of c, at time c is

\[= (c,) \cos(nt +=   ),\]
(2) where (c,), o denote the initial values of the
quantities affected with the suffix s. Now in the problem in hand (c,),= and
(c,), is determined by

\[n b= 7 =7 Y \sin()= 0 (8),\]
if Y denote the force with which the string is
held aside at the point b. Hence at time c

\[c =  \sin(5) - \cos(N)\]

(4), and by (1) of Article 128

\[\cos(nt) y= Ip  \sin(T) \sin(T)\]

(5), where n = g : 1.

The symmetry of the expression (5) in and 6 is an example of the principle of
Article 107.

The problem of determining the subsequent motion of a string set into vibration
by an impulse acting at the point b, may be treated in a similar manner.
Integrating (6) of Article 128 over the duration of the impulse, we find
ultimately, with the same notation as before,

Dangling\footnote{A like result ensues when the point which is damped is at the
same distance from one end of the string as the point of excitation is from the
other end.}

if Y' dt be denoted by Y,, At the same time (c,), = 0, 80 that by

(2) at time c

\[\sin(nt) \sin() \sin(nr).\]

(6). The series of component vibrations is less convergent for a struck than
for a plucked string, as the preceding expressions shew. The reason is that in
the latter case the initial value of y is continuous, and only ay
discontinuous, while in the former it is

y itself that makes a sudden spring. See Article 32, 101.

The problem of a string set in motion by an impulse may also be solved by the
general formulae (7) and (8) of Article 128. The force finds the string at rest
at c=0, and acts for an infinitely short time from t=0 to t=7. Thus (c,), and
(c,), vanish, and (7) of Article 128 reduces to. :

\[c = Ga \sin(ne) P de,\]

while by (8) of Article 128

\[\sin(Yar) = \sin(Y),\]

Hence, as before, =F Hitherto we have supposed the disturbing force to be
concentrated at a single point. If it be distributed over a distance 8 on
either side of b, we have only to integrate the expressions (6) and (7) with
respect to 6, substituting, for example, in (7) in place of Y,

\[= Y \sin()=- SID NE. ..(Ze\]

\[\sin(o), \sin(Fa),\]

If Y be constant between the limits, this reduces to

\[Y \sin(7)\]

(8). The principal effect of the distribution of the force is to render the
series for y more convergent.


130. The problem which will next engage our attention is that of the pianoforte
wire. The cause of the vibration is here the blow of a hammer, which is
projected against the string, and after the impact rebounds. But we should not
be justified in assuming, as in the last section, that the mutual action
occupies so short a time that its duration may be neglected. Measured by the
standards of ordinary life the duration of the contact is indeed very small,
but here the proper comparison is with the natural periods of the string. Now
the hammers used to strike the wires of a pianoforte are covered with several
layers of cloth for the express purpose of making them more yielding, with the
effect of prolonging the contact. The rigorous treatment of the problem would
be difficult, and the solution, when obtained, probably too complicated to be
of use; but by introducing a certain simplification Helmholtz has obtained a
solution representing all the essential features of the case. He remarks that
since the actual yielding of the string must be slight in comparison with that
of the covering of the hammer, the law of the force called into play during the
contact must be nearly the same as if the string were, absolutely fixed, in
which case the force would vary very nearly as a circular function. We shall
therefore suppose that at the time t= 0, when there are neither velocities nor
displacements, a force \(F \sin(pt)\) begins to act on the string at 2=6, and
continues through half a period of the circular function, that is, until t= 1+
p, after which the string is once more free. The magnitude of p will depend on
the mass and elasticity of the hammer, but not to any great extent on the
velocity with which it strikes the string.

The required solution is at once obtained by substituting for , in the general
formula (7) of Article 128 its value given by

\[= F \sin(7) a i n pt (1),\]
the range of the integration being from 0 to We
find (5)

\[6,=ja \sin(p) f \sin(n) (c-t) \sin(pt)de 4p \cos(5) and , = Fin \sin(a) (t- 3) ,\]
(2) and the final solution for y becomes, if we substitute for x and p their values,

(3) We see that all components vanish which have a node at the point of excitement,
but this conclusion does not depend on any particular law of force. The
interest of the present solution lies in the information that may be elicited
from it as to the dependence of the resulting vibrations on the duration of
contact. If we denote the ratio of this quantity to the fundamental period of
the string by v, so that v = 7a: 2pl, the expression for the amplitude of the
component s is

\[\cos() \cos(s) s (1 - 48 ) 4s v  \sin(T)\]

(4). We fall back on the case of an impulse by putting v=0 and

When is finite, those components disappear, whose periods are , , , of the
duration of contact; and when s is very great, the series converges with s .
Some allowance must also be made for the finite breadth of the hammer, the
effect of which will also be to favour the convergence of the series.

The laws of the vibration of strings may be verified, at least in their main
features, by optical methods of observation - either, with the
vibration-microscope, or by a tracing point recording the character of the
vibration on a: revolving drum. This character depends on two things, the mode
of excitement, and the point whose motion is selected for observation. Those
components do not appear which have nodes either at the point of excitement, or
at the point of observation. The former are not generated, and the latter do
not manifest themselves. Thus the simplest motion is obtained by plucking the
string at the centre, and observing one of the points of trisection, or vice
versa. In this case the first harmonic which contaminates the purity of the
principal vibration is the fifth component, whose intensity is usually
insufficient to produce much disturbance.

page 191, pdf page 205 192
dynamical theory of -the vibration of strings may be employed to test the laws
of hearing, and the necessary experiments are easily carried out upon a grand
pianoforte. Having freed a string, say c, from its damper by pressing the
digital, pluck it at one-third of its length. According to Young s theorem the
third component vibration is not excited then, and in correspondence with that
fact the ear fails to detect the component A slight displacement of the point
plucked brings in again; and if a resonator (jgf) be useful to assist the ear,
it is only with difficulty that the point can be hit with such precision as
entirely to extinguish the tone. Experiments of this kind shew that the ear
analyses the sound of a string into precisely the same constituents as are
found by sympathetic resonance, that is, into simple tones, according to Ohm’s
definition of this conception.  Such experiments are also well adapted to shew
that it is not a mere play of imagination when we hear overtones, as some
people believe it is on hearing them for the first time

If, after the string has been sounded loudly by striking the digital, it be
touched with the finger at one of the points of trisection, all components are
stopped except the 3rd, 6th, etc., so that these are left isolated. The
inexperienced observer is usually surprised by the loudness of the residual
sound, and begins to appreciate the large part played by overtones.]

131. The case of a periodic force is included in the general solution of
Article 128, but we prefer to follow a somewhat different method, in order to
make an extension in another direction. We have hitherto taken no account of
dissipative forces, but we will now suppose that the motion of each element of
the string is resisted by a force proportional to its velocity. The partial
differential equation becomes

FY 4 i dy a FY He Y (1), by means of which the subject may be treated.  But it
is still simpler to avail ourselves of the results of the last chapter,
remarking that in the present case the friction-function F is of the same form
as T. In fact

Bem fad. BF P (2),

=1 where c,, ,,... are the normal coordinates, by means of which fT and V are
reduced to sums of squares. The equations of motion are therefore simply

d, + xd, + n, = Do (3),

of the same form as obtains for systems with but one degree of freedom.  It is
only necessary to add to what was said in Chapter III., that since is
independent of s, the natural vibrations subside in such a manner that the
amplitudes maintain their relative values.

If a periodic force \(F \cos(pt)\) act at a single point, we have

\[=F \sin(s)  \cos(P)\]

(4), and Article 46 .

\[c- \sin(oP 08 (pt -) (5),\]

where

\[tan e= top (6).\]

If among the natural vibrations there be any one nearly isochronous with
\(\cos(pt)\), then a large vibration of that type will be forced, unless indeed
the point of excitement should happen to fall near a node. In the case of exact
coincidence, the component vibration in question vanishes; for no force applied
at a node can generate it, under the present law of friction, which however, it
may be remarked, is very special in character. If there be no friction, = 0,
and

\[= ao \sin()= y\]

(7), which would make the vibration infinite, in the case of perfect
isochronism, unless \(\sin(oe) =0\).

The value of y is here, as wa

\[y=49, \sin()= T a+ cg, \sin(277) r+ c, \sin(r) (8).\]

132. The preceding solution is an example of the use of normal coordinates in a
problem of forced vibrations. It is of course to free vibrations that they are
more especially applicable, and they may generally be used with advantage
throughout, whenever the system after the operation of various forces is
ultimately left to itself. Of this application we have already had examples.

In the case of vibrations due to periodic forces, one advantage of the use of
normal coordinates is the facility of comparison with the equilibrium theory,
which it will be remembered is the theory of the motion on the supposition that
the inertia of the system may be left out of account. If the value of the
normal coordinate c, on the equilibrium theory be A, \(\cos(pt)\), then the
actual value will be given by the equation

\[,= WA, rs \cos( \Phi ). (1),\]

so that, when the result of the equilibrium theory is known and can readily be
expressed in terms of the normal coordinates, the true solution with the
effects of inertia included can at once be written down.

In the present instance, if a force

\[F' \cos(pt)\]

of very long period act at the point b of the string, the result of the
equilibrium theory, in accordance with which the string would at any moment
consist of two straight portions, will be

\[pd, = or \sin(me) \cos(pt) (2),\]


from which the actual result for all values of p is derived by simply writing
\(n - p\) in place of n'.

The value of y in this and similar cases may however be expressed in finite
terms, and the difficulty of obtaining the finite expression is usually no
greater than that of finding the form of the normal functions when the system
is free. Thus in the equation of motion

=a ty

suppose that Y varies as \(\cos(mat)\) . The forced vibration will then satisfy
. 1; ay + my = - fe Cee (8).

If Y=0, the investigation of the normal functions requires the solution of

and a subsequent determination of m to suit the boundary conditions. In the
problem of forced vibrations m is given, and we have only to supplement any
particular solution of (8) with the complementary function containing two
arbitrary constants. This function, apart from the value of m and the ratio of
the constants, is of the same form as the normal functions; and all that
remains to be effected is the determination of the two constants in accordance
with the prescribed boundary conditions which the complete solution must
satisfy. Similar considerations apply in the case of any continuous system.

183. If a periodic force be applied at a single point, there are two distinct
problems to be considered; the first, when at the point a =, a given periodic
force acts; the second, when it is the actual motion of the point 6 that is
obligatory. But it will be convenient to treat them together.

The usual differential mee

\[+eQ HOT \cos(e) (1),\]

is satisfied over both the parts into which the string is divided at b, but is
violated in crossing from one to the other.

In order to allow for a change in the arbitrary constants, we must therefore
assume distinct expressions for y, and afterwards introduce the two conditions
which must be satisfied at the point of junction.  These are

(1) That there is no discontinuous change in the value of y;

(2) That the resultant of the tensions acting at 6 balances the impressed
force. ,

Thus, if
\[F' \cos(pt)\]
be the force, the second condition gives
\[+eQ HOT \cos(e) (1),\]
\[TA (2) 4+ F \cos(pt) = 0  (2),\]
where A (4) denotes the alteration in the value of dy incurred dz; in crossing
the point x= 6 in the positive direction.

We shall, however, find it advantageous to replace \(\cos(pt)\) by the complex
exponential e , and finally discard the imaginary part, when the symbolical
solution is completed. On the assumption that y varies as e, the differential
equation becomes

where X is the complex constant, (4).

The most general solution of (3) consists of two terms, proportional
respectively to \(\sin(Aw)\), and \(\cos(Aw)\); but the condition to be
satisfied at = 0, shews that the second does not occur here. Hence if ye be the
value of y at =6,

\[\sin(Ax)  \sin(d) dx (5),\]
is the solution applying to the first part of the
string from =0 toz=b, In like manner it is evident that for the second part we
shall have

\[\sin(a) Y= b.\]

(6). If y be given, these equations constitute the symbolical solution of the
problem; but if it be the force that be given, we require further to know the
relation between it and y.


Differentiation of (5) and (6) and substitution in the equation analogous to
(2) gives

\[F \sin(Xd) \sin(A(1 - 5))\]

\[YT, X \sin(Al V).\]

Thus

\[F \sin(dx \sin(d(l-b) Ip \sin(Al)\]
 from =0 to a=
\[\sin(A(J- 2) \sin(ND)\]
in
\[X \sin(Al)\]

(8) These equations exemplify the general law of reciprocity proved in the last
chapter; for it appears that the motion at x due to the force at (missing) is
the same as would have been found at 6, had the force acted at x.

In discussing the solution we will take first the case in which there is no
friction. The coefficient is then zero; while is real, and equal to p+a. The
real part of the solution, corresponding to the force \(F' \cos(pt)\), is found
by simply putting \(\cos(pt)\) for e in (8), but it seems scarcely necessary to
write the equations again for the sake of so small a change. The same remark
applies to the forced motion given in terms of .

It appears that the motion becomes infinite in case the force is isochronous
with one of the natural vibrations of the entire string, unless the point of
application be a node; but in practice it is not easy to arrange that a string
shall be subject to a force of given magnitude. Perhaps the best method would
be to attach a small mass of iron, attracted periodically by an electromagnet,
whose coils are traversed by an intermittent current. But unless some means of
compensation were devised, the mass would have to be very small in order to
avoid its inertia introducing a new complication.

A better approximation may be obtained to the imposition of an obligatory
motion. A massive fork of low pitch, excited by a bow or sustained in permanent
operation by electromagnetism, executes its vibration in approximate
independence of the reactions of any light bodies which may be connected with
it. In order therefore to subject any point of a string to an obligatory
transverse motion, it is only necessary to attach it to the extremity of one
prong of such a fork, whose plane of vibration is perpendicular to the length
of the string. This method of exhibiting the forced vibrations of a string
appears to have been first used by Melde.

Donkin's Acoustics, p. 121.

Another arrangement, better adapted for aural observation, has been employed by
Helmholtz. The end of the stalk of a powerful tuning-fork, set into vibration
with a bow, or otherwise, is pressed against the string. It is advisable to
file the surface, which comes into contact with the string, into a suitable
(saddle shaped) form, the better to prevent slipping and jarring.

Referring to (5) we see that, if \(\sin(Xb)\) vanished, the motion (according
to this equation) would become infinite, which may be taken to prove that in
the case contemplated, the motion would really become great, so great that
corrections, previously insignificant, rise into importance. Now \(\sin(Ab)\)
vanishes, when the force is isochronous with one of the natural vibrations of
the first part of the string, supposed to be held fixed at 0 and b.

When a fork is placed on the string of a monochord, or other instrument
properly provided with a sound-board, it is easy to find by trial the places of
maximum resonance. A very slight displacement on either side entails a
considerable falling off in the volume of the sound. The points thus determined
divide the string into a number of equal parts, of such length that the natural
note of any one of them (when fixed at both ends) is the same as the note of
the fork, as may readily be verified. The important applications of resonance
which Helmholtz has made to purify a simple tone from extraneous accompaniment
will occupy our attention later,

134. Returning now to the general case where A is complex, we have to extract
the real parts from (5), (6), (8) of Article 183. For this purpose the sines
which occur as factors, must be reduced to the form Re. Thus let
\[\sin(Ma = Ry eM (1),\]
with a like notation for the others. From (5) Article 138 we shall thus obtain
\[Y= 1A \cos(pt) + y-p (2),\]

from x=0 to w=6,

and from (6) Article 133

\[y=7 R. \cos(pt)\]

bd from x=b to x=l,

corresponding to the obligatory motion \(y =y \cos(pt)\) at bd. By a similar
process from (8) Article 133, if

we should obtain from

corresponding to the impressed force \(F' \cos(pt)\) at b. It remains to obtain
the forms of ,, e,, etc. The values of and 8 are determined by (4), and

\[= \sin(Aw) = \sin(ax) \cos(tBx)-+ \cos(ax) B \sin(ie)\]

\[= \sin(ax) -5 +i \cos(az)\]

so that

\[BR = \sin(ax) (4e) + \cos(ax) a wee\]

(6),

\[tan e, =\]

(7);

while (8). This completes the solution. If the friction be very small, the
expressions may be simplified. For instance, in this case, to a sufficient
approximation,
\[1 P -  a, \pi P a= B=-55) V+ B=)\]
so that corresponding to
the obligatory motion at \(b y=+y \cos(pt)\) the amplitude of the motion
between 2 = 0 and x = b is, approximately

\[\sin(t PE 4 ex \cos(t)\]

\[\sin(Fa) cos\]
which becomes great, but not infinite, when \(\sin(0)\), or
the point of application is a node.

If the imposed force, or motion, be not expressed by a single harmonic term, it
must first be resolved into such. The preceding solution may then be applied to
each component separately, and the results added together. The extension to the
case of more than one point of application of the impressed forces is also
obvious. To obtain the most general solution satisfying the conditions, the
expression for the natural vibrations must also be added; but these become
reduced to insignificance after the motion has been in progress for a
sufficient time.

The law of friction assumed in the preceding investigation is the only one
whose results can be easily followed deductively, and it is sufficient to give
a general idea of the effects of dissipative forces on the motion of a string.
But in other respects the conclusions drawn from it possess a fictitious
simplicity, depending on the fact that the friction function is similar in form
to 7, which makes the normal coordinates independent of each other. In almost
any other case (for example, when but a single point of the string is retarded
by friction) there are no normal coordinates properly so called. There exist
indeed elementary types of vibration into which the motion may be resolved, and
which are perfectly independent, but these are essentially different in
character from those with which we have been concerned hitherto, for the
various parts of the system (as affected by one elementary vibration) are not
simultaneously in the same phase. Special cases excepted, no linear
transformation of the coordinates (with real coefficients) can reduce 7', F,
and V together to a sum of squares.

If we suppose that the string has no inertia, so that 7'=0, F and V may then be
reduced to sums of squares. This problem is of no acoustical importance, but it
is interesting as being mathematically analogous to that of the conduction and
radiation of heat in a bar whose ends are maintained at a constant temperature,


135. Thus far we have supposed that at two fixed points, x=0 and =/ , the
string is held at rest. Since absolute fixity cannot be attained in practice,
it is not without interest to inquire in what manner the vibrations of a string
are liable to be modified by a yielding of the points of attachment; and the
problem will furnish occasion for one or two remarks of importance. For the
sake of simplicity we shall suppose that the system is symmetrical with
reference to the centre of the string, and that each extremity is attached to a
mass M (treated as unextended in space), and is urged by a spring (4) towards
the position of equilibrium. If no frictional forces act, the motion is
necessarily resolvable into normal vibrations. Assume

\[y = 2 \sin(mx) + B \cos(mz) \cos(mat - )\]

(1). The conditions at the ends are that

when w=0, My+py= 7,

when aw=1, My j + ny =- 7,9 which give

two equations, sufficient to determine m, and the ratio of 8 to a.  Eliminating
the latter ratio, we find

\[2y \tan(ml) = -p2 ever\]

(4), Matin if for brevity we write v for

Equation (3) has an infinite number of roots, which may be found by writing \(
\tan() \) for v,so that \(\tan(ml) = \tan(20) \), and the result of adding
together all the corresponding particular solutions, each with its two
arbitrary constants a and c, is necessarily the most general solution of which
the problem is capable, and is therefore adequate to represent the motion due
to an arbitrary initial distribution of displacement and velocity. We infer
that any function of a may be expanded between x=0 and =1 in a series of terms

\[c, (v, \sin(m x) + \cos(mx) + , (v, \sin(mx) + \cos(mr) +...4.(5),\]
 m,, m,,
etc. being the roots of (3) and ,, v,, etc. the corresponding values of v. The
quantities c,, c,, are the normal coordinates of the system.

From the symmetry of the system it follows that in each normal vibration the
value of y is numerically the same at points equally distant from the middle of
the string, for example, at the two ends, where x=0 and x=/. Hence p, \(\sin(m)
+ \cos(ml) =+1,\) as may be proved also from (4).

The kinetic energy 7 of the whole motion is made up of the energy of the
string, and that of the masses M. Thus to;

\[T= al   (v \sin(mx) + \cos(mz)  dx\]

\[+4M (6, 46,4.) +4 (9, \sin(ml) + \cos(m,) +\]

But by the characteristic property of normal coordinates, terms containing
their products cannot be really present in the expression for 7, so that

\[p f (v, \sin(m) + \cos(m,z) (v, \sin(mx) + \cos(mx) dx\]
\[+M+M (v \sin(m,1)+ \cos(m) \sin(m,l) + \cos(m,l) = 0 (6),\]
if r and s be different.

This theorem suggests how to determine the arbitrary constants, so that the
series (5) may represent an arbitrary function y. Take the expression .

t

\[p   y (v \sin(m,2-+ \cos(m,x)dx+My,+ My,(v, \sin(m,1 + \cos(m, )...(7),\]

0

and substitute in it the series (5) expressing y. The result is a series of
terms of the type

\[Pp   iP (v, \sin(me + \cos(m,2) (v, \sin(m,a + \cos(m,x) dz\]

\[+ Md + Me (v \sin(m_1 + \cos(m ) (v \sin(m_1 + \cos(m)),\]

all of which vanish by (6), except the one for which r=s, Hence c, is equal to
the expression (7), divided by t

\[Pp i  \sin(m,x + \cos(m,x)" dx + M+ M (v \sin(m1 + \cos(m, ) ...(8),\]

and thus the coefficients of the series are determined. If M=0, even although y
be finite, the process is of course much simpler, but the unrestricted problem
is instructive. So much stress is often laid on special proofs of Fourier's and
Laplace's series, that the student is apt to acquire too contracted a view of
the nature of those important results of analysis.


We shall now shew how Fourier's theorem in its general form can be deduced from
our present investigation. Let M=0; then if = 0, the ends of the string are
fast, and the equation determining m becomes \( \tan(ml) =0\), or ml = sr, as
we know it must be. In this case the series for y becomes

me ome.

\[A \sin(y)= A, \sin(i) +4 \sin()+ A \sin()= + (9),\]
which must be general
enough to represent any arbitrary functions of 2, vanishing at 0 and J, between
those limits. But now suppose that yu is zero, M still vanishing. The ends of
the string may be supposed capable of sliding on two smooth rails perpendicular
to its length, and the terminal condition is the vanishing of wy. The equation
in m is the same as before; and we learn that any function y' whose rates of
variation vanish at 2=0 and a=/, can be expanded in a series

\[y = B \cos() B \cos() = + B \cos(2) 4 teases (10).\]

This series remains unaffected when the sign of x is changed, and the first
series merely changes sign without altering its numerical magnitude. If
therefore y' be an even function of 2, (10) represents it from -/ to +7. And in
the same way, if y be an odd function of z, (9) represents it between the same
limits.

Now, whatever function of (a) may be, it can be divided into two parts, one of
which is even, and the other odd, thus:

so that, if c (2) be such that c (-1) = and c'(-) =c' , it can be represented
between the limits + 7 by the mixed series

\[+ A \sin(72) + B \cos(4)\]

(11).

\[A, \sin(7) +B, cos\]

This series is periodic, with the period 22. If therefore c (z) possess the
same property, no matter what in other respects its character may be, the
series is its complete equivalent. This is Fourier's theorem.

We now proceed to examine the effects of a slight yielding of the supports, in
the case of a string whose ends are approximately fixed.  The quantity v may be
great, either through or through M. We shall confine ourselves to the two
principal cases, (1) when p is great and M vanishes, (2) when m vanishes and M
is great.

In the first case y= Tm'

and the equation in m is approximately \(\tan(ml) = -2\).

Assume ml =s1 + 2, where w is small; then

e=27 - 6m approximately,

and m= er (1 - =) se (12).

Dangling\footnote{ The best `system' for proving Fourier's theorem from dynamical
considerations is an endless chain stretched round a smooth cylinder (Article 139),
or a thin re-entrant column of air enclosed in a ring-shaped tube.}

To this order of approximation the tones do not cease to form a harmonic scale,
but the pitch of the whole is slightly lowered. The effect of the yielding is
in fact the same as that of an increase in the length of the string in the
ratio 1:1 42h as might have been anticipated. The result is otherwise if
vanish, while M is great. Here y=- Ma m an a and \(\tan(ml) = 21\),
approximately. Mam PP y: Hence 271 ml = sr + 53 - M ism cee .(13). The effect
is thus equivalent to a decrease in J in the ratio . 27,1 til ar ia and
consequently there is a rise in pitch, the rise being the greater the lower the
component tone. It might be thought that any kind of yielding would depress the
pitch of the string, but the preceding investigation shews that this is not the
case. Whether the pitch will be raised or lowered, depends on the sign of vy,
and this again depends on whether the natural note of the mass M urged by the
spring yw is lower or higher than that of the component vibration in question.

135. . FINITE LOAD, 161

136. The problem of an otherwise uniform string carrying a finite load M at =b
can be solved by the formulae investigated in Article 133. For, if the force
\(F' \cos(pt)\) be due to the reaction against acceleration of the mass M,

which combined with equation (7) of Article 133 gives, to determine the
possible values of (or p : a),

\[a' M \sin(A) \sin(d (1 - b)) = T \sin(Al) (2).\]

The value of y for any normal vibration corresponding to A is

\[y = P \sin(dx \sin(A) (J- b) \cos(a) dt - e\]

from 2=0 t=b

\[y= P \sin(d(t-a) \sin(Xb) \cos(at-e) (8)\]

from =b to x=1

where P and c are arbitrary constants.

It does not require analysis to prove that any normal components which have a
node at the point of attachment are unaffected by the presence of the load. For
instance, if a string be weighted at the centre, its component vibrations of
even orders remain unchanged, while all the odd components are depressed in
pitch. Advantage may sometimes be taken of this effect of a load, when it is
desired for any purpose to disturb the harmonic relation of the component
tones.

If M be very great, the gravest component is widely separated in pitch from all
the others. We will take the case when the load is at the centre, so that
b=b/(1-b). The equation in

then becomes

where pl : M, denoting the ratio of the masses of the string and the load, is a
small quantity which may be called a. The first

(4),

root corresponding to the tone of lowest pitch occurs when 421 is small, and
such that

nearly,

and the periodic time is given by

(5).

The second term constitutes a correction to the rough value obtained in a
previous chapter (Article 52), by neglecting the inertia of the string
altogether. That it would be additive might have been expected, and indeed the
formula as it stands may be obtained from the consideration that in the actual
vibration the two parts of the string are nearly straight, and may be assumed
to be exactly so in computing the kinetic and potential energies, without
entailing any appreciable error in the calculated period. On this supposition
the retention of the inertia of the string increases the kinetic energy
corresponding to a given velocity of the load in the ratio of M: M+ pl, which
leads to the above result. This method has indeed the advantage in one respect,
as it might be applied when p is not uniform, or nearly uniform. All that is
necessary is that the load M should be sufficiently predominant.

whence

Fig. 21.

There is no other root of (4), until \(\sin(A)=0\), which gives the second
component of the string-a vibration independent of the load. The roots after
the first occur in closely contiguous pairs; for one set is given by 4Al=e7,
and the other approximately by l =s7 +P in which the second term is small.

The two types of vibration for s = 1 are shewn in the figure.

The general formula (2) may also be applied to find the effect of a small load
on the pitch of the various components,

137. Actual strings and wires are not perfectly flexible. They oppose a certain
resistance to bending, which may be divided into two parts, producing two
distinct effects. The first is called viscosity, and shews itself by damping
the vibrations. This part produces no sensible effect on the periods. The
second is conservative in its character, and contributes to the potential
energy of the system, with the effect of shortening the periods. A complete
investigation cannot conveniently be given here, but the case which is most
interesting in.its application to musical instruments, admits of a sufficiently
simple treatment,

When rigidity is taken into account, something more must be specified with
respect to the terminal conditions than that y vanishes. Two cases may be
particularly noted :

(1) When the ends are clamped, so that oy = 0 at the ends.

(2) When the terminal directions are perfectly free, in which case ay = 0, dx
It is the latter which we propose now to consider. If there were no rigidity,
the type of vibration would be y sin, satisfying the second condition.

The effect of the rigidity might be slightly to disturb the type; but whether
such a result occur or not, the period calculated from the potential and
kinetic energies on the supposition that the type remains unaltered is
necessarily correct as far as the first order of small quantities (Article 88).

Now the potential energy due to the stiffness is expressed by

where is a quantity depending on the nature of the material and on the form of
the section in a manner that we are not now prepared to examine.  The form of
8V is evident, because the force required to bend any element ds is
proportional to ds, and to the amount of bending already effected, that is to
ds+p. The whole work which must be done to produce a curvature 1+p in ds is
therefore proportional to ds+p; while to the approximation to

which we work ds=da, and

Thus, if \(y=g \sin()=,\)

and the period of c is given by

(2),

If r denote what the period would become if the string were endowed with
perfect flexibility. It appears that the effect of the stiffness increases
rapidly with the order of the component vibrations, which cease.to belong to a
harmonic scale. However, in the strings employed in music, the tension is
usually sufficient to reduce the influence of rigidity to insignificance.

The method of this section cannot be applied without modification to the other
case of terminal condition, namely, when the ends are clamped. In their
immediate neighbourhood the type of vibration must differ from that assumed by
a perfectly flexible string by a quantity, which is no longer small, and whose
square therefore cannot be neglected. We shall return to this subject, when
treating of the transverse vibrations of rods.

138. There is one problem relating to the vibrations of strings which we have
not yet considered, but which is of some practical interest, namely, the
character of the motion of a violin (or cello) string under the action of the
bow. In this problem the modus operandi of the bow is not sufficiently
understood to allow us to follow exclusively the \textit{a priori} method: the
indications of theory must be supplemented by special: observation. By a
dexterous combination of evidence drawn from both sources Helmholtz has
succeeded in determining the principal features of the case, but some of the
details are still obscure.

Since the note of a good instrument, well handled, is musical, we infer that
the vibrations are strictly periodic, or at least that strict periodicity is
the ideal. Moreover-and this is very important-the note elicited by the bow has
nearly, or quite, the same pitch as the natural note of the string. The
vibrations, although forced, are thus in some sense free. They are wholly
dependent for their maintenance on the energy drawn from the bow, and yet the
bow does not determine, or even sensibly modify, their periods. We are reminded
of the self-acting electrical interrupter, whose motion is indeed forced in the
technical sense, but has that kind of freedom which consists in determining
(wholly, or in part) under what influences it shall come.

But it does not at once follow from the fact that the string vibrates with its
natural periods, that it conforms to its natural types. If the coefficients of
the Fourier expansion

\[. 2; y=  \sin(F) r+  \sin()+ \sec(a)\]

be taken as the independent coordinates by which the configuration of the
system is at any moment defined, we know that when there is no friction, or
friction such that F' oc 7', the natural vibrations are expressed by making
each coordinate a simple harmonic (or quasi-harmonic) function of the time;
while, for all that has hitherto appeared to the contrary, each coordinate in
the present case might be any function of the time periodic in time r. But a
little examination will shew that the vibrations must be sensibly natural in
their types as well as in their periods.

The force exercised by the bow at its point of application may be expressed by

Y= 3A () = -); so that the equation of motion for the coordinate c, is

\[b + R = Z \sin(5)  \cos(),\]

b being the point of application. Each of the component parts of, will give a
corresponding term of its own period in the solution, but the one whose period
is the same as the natural period of c, will rise enormously in relative
importance. Practically then, if the damping be small, we need only retain that
part of c which depends on A () =- ), that is to say, we may regard the
vibrations as natural in their types,

Another material fact, supported by evidence drawn both from theory and aural
observation,.is this, All component vibrations are absent which have a node at
the point of excitation. In order, however, to extinguish these tones, it is
necessary that the coincidence of the point of application of the bow with the
node should be very exact. A very small deviation reproduces the missing tones
with considerable strength.''

The remainder of the evidence on which Helmholtz' theory rests, was derived
from direct observation with the vibration microscope. As explained in Chapter
11., this instrument affords a view of the curve representing the motion of the
point under observation, as it would be seen traced on the surface of a
transparent cylinder. In order to deduce the representative curve in its
ordinary form, the imaginary cylinder must be conceived to be unrolled, or
developed, into a plane.

The simplest results are obtained when the bow is applied at a node of one of
the higher components, and the point observed is one of the other nodes of the
same system. If the bow works fairly so as to draw out the fundamental tone
clearly and strongly, the representative curve is that shewn in figure 22;
where the abscissa correspond to the time (AB being a complete period), and the
ordinates represent the displacement. The remarkable fact is disclosed that the
whole period + may be divided into two parts r, and r-7,, during each of which
the velocity of the observed point is constant; but the velocities to and fro
are in general unequal.

Fig. 22.

We have now to represent this curve by a series of harmonic terms. If the
origin of time correspond to the point 4, and

Dangling\footnote{ Donkin's Acoustics, p. 181.}

138. VIOLIN STRING. 167

AD=FU=y, Fourier's theorem gives; sin - (y= mt, (7 -7,) 2 FF T With respect to
the value of we know that all those com T" point of observation), because under
the circumstances of the case the bow cannot generate them. There is therefore
reason to suppose that r,:7r=2,:1; and in fact observation proves that AC: CB
(in the figure) is equal to the ratio of the two parts into which the string is
divided by the point of observation. Now the free vibrations of the string are
represented in general by

y n (7) 4 (Ts B, ()=

and this at the point =x, must agree with (1). For convenience of comparison,
we may write A (2) 4 B, (Pert C () + D (2) en (- i), e T DA and it then appears
that C, = 0.

We find also to determine D, (a) = T whence 2y7 1 find rr T, (r- 7.) 3.  C .
(2), unless ()=0,

In the case reserved, the comparison leaves D, undetermined, but we know on
other grounds that D, then vanishes. However, for the sake of simplicity, we
shall suppose for the present that D, is always given by (2). If the point of
application of the bow do not coincide with a node of any of the lower
components, the error committed will be of no great consequence.

On this understanding the complete solution of the problem is

By (2)

The amplitudes of the components are therefore proportional to s, In the case
of a plucked string we found for the corresponding function sin ox , which is
somewhat similar. If the string be plucked at the middle, the even components
vanish, but the odd ones follow the same law as obtains for a violin string.
The equation (3) indicates that the string is always in the form of two
straight lines meeting at an angle. In order more conveniently to shew this,
let us change the origin of the time, and the constant multiplier, so that srx
. srt \(= \Psi\) psn SID (A), will be the equation expressing the form of the string
at any time.

Now we know (Article 127) that the equation of the pair of lines proceeding
from the fixed ends of the string, and meeting at a point whose coordinates are
a, 8, is

2AP 31 (om gin ST

ma (l-a) 23 rn Thus at the time c, (6) represents.such a pair of lines, meeting
at the point whose coordinates are given by

Ar a(l-a) - = 48, (a) i n 2278 Sao se sin

These equations indicate that the projection on the axis of x of the point of
intersection moves uniformly backwards and forwards between 2=0 and =1, and
that the point of intersection itself is situated on one or other of two
parabolic arcs, of which the equilibrium position of the string i is a common
chord.

Since the motion of the string as thus defined by that of the point of
intersection of its two straight parts, has no especial relation to a, (the
point of observation), it follows that, according to these equations, the same
kind of motion might be observed at any other point.  And this is approximately
true. But the theoretical result, it will be remembered, was only obtained by
assuming the presence in certain proportions of component vibrations having
nodes at 2, though in fact their absence is required by mechanical laws. The
presence or absence of these components is a matter of indifference when a node
is the point of observation, but not in any other case. When the node is
departed from, the vibration curve shews a series of ripples, due to the
absence of the components in question. Some further details will be found in
Helmholtz and Donkin.

138. STRINGS STRETCHED ON CURVED SURFACES. 169

The sustaining power of the bow depends upon the fact that solid friction is
less at moderate than at small velocities, so that when the part of the string
acted upon is moving with the bow (not improbably at the same velocity), the
mutual action is greater than when the string is moving in the opposite
direction with a greater relative velocity. The accelerating effect in the
first part of the motion is thus not entirely neutralised by the subsequent
retardation, and an outstanding acceleration remains capable of maintaining the
vibration in spite of other losses of energy. A curious effect of the same
peculiarity of solid friction has been observed by Mr Froude, who found that
the vibrations of a pendulum swinging from a shaft might be maintained or even
increased by causing the shaft to rotate.

[Another case in which the vibrations of a string are maintained is that of the
Aeolian Harp. It has often been suggested that the action of the wind is
analogous to that of a bow; but the analogy is disproved by the observation
that the vibrations are executed in a plane transverse to the direction of the
wind. The true explanation involves hydrodynamical theory not yet developed.]

139. A string stretched on a smooth curved surface will in equilibrium lie
along a geodesic line, and, subject to certain conditions of stability, will
vibrate about this configuration, if displaced. The simplest case that can be
proposed is when the surface is a cylinder of any form, and the equilibrium
position of the string is perpendicular to the generating lines. The student
will easily prove that the motion is independent of the curvature of the
cylinder, and that the vibrations are in all essential respects the same as if
the surface were developed into a plane. The case of an endless string, forming
a necklace round the cylinder, is worthy of notice.

In order to illustrate the characteristic features of this class of problems,
we will take the comparatively simple example of a string stretched on the
surface of a smooth sphere, and lying, when in equilibrium, along a great
circle. The coordinates to which it will be most convenient to refer the system
ate the latitude measured from the great circle as equator, and the longitude c
measured along it. If the radius of the sphere be a, we have;

The extension of the string is denoted by

Now

so that

approximately. Thus

If the ends be fixed, "

and the equation of virtual velocities is

whence, since is arbitrary;

(3)

This is the equation of motion.

If we assume 0 (pt), we get

(4),

of which the solution, subject to the condition that vanishes with c, is

6=A (Fr p oly P (ME (5). The remaining condition to be satisfied is that 6
vanishes when ap =I, or g=a, if a=l+a, This gives T, min T, min . 1 7 1 - = --+
p ap (ai 1) p (Z =) (6), where m is an integer. 1 Cambridge Mathematical Tripos
Examination, 1876.

139. VARIABLE DENSITY. 171

The normal functions are thus of the same form as for a straight string, viz.

6=A (P C08 pt (7),

but the series of periods is different. The effect of the curvature is to make
each tone graver than the corresponding tone of a straight string. If a, one at
least of the values of p' is negative, indicating that the corresponding modes
are unstable. If a=, p, is zero, the string being of the same length in the
displaced position, as when 6 = 0.

A similar method might be applied to calculate the motion of a string stretched
round the equator of any surface of revolution.

[For a more general treatment of this question see Michell, Messenger of
Mathematics, vol. XIX. p. 87, 1890.]

140. The approximate solution of the problem for a vibrating string of nearly
but not quite uniform longitudinal density has been fully considered in Chapter
IV. Article 91, as a convenient example of the general theory of approximately
simple systems. It will be sufficient here to repeat the result. If the density
be p,+ dp, the period +, of the r component vibration is given by

(1).

If the irregularity take the form of a small load of mass m at the point =, the
formula may be written om pel + a i n int (2).

These values of + are correct as far as the first power of the small quantities
5p and m, and give the means of calculating a correction for such slight
departures from uniformity as must always occur in practice.

As might be expected, the effect of a small load vanishes at nodes, and rises
to a maximum at the points midway between consecutive nodes. When it is desired
merely to make a rough estimate of the effective density of a nearly uniform
string, the formula indicates that attention is to be given to the
neighbourhood of loops rather than to that of nodes.

[The effect of a small variation of density upon the period is the same whether
it occur at a distance x from one end of the string, or at an equal distance
from the other end. The mean variation at points equidistant from the centre is
all that we need regard, and thus no generality will be lost if we suppose that
the density remains symmetrically distributed with respect to the centre.  Thus
we may write d + 1 (3) where
2f So - 2 r (4).
In this equation Sp may be expanded from 0 to 2 in the sense
- = .4, + di cos + ... + ri cos + ...  P t (5).  where ().  (7).
Accordingly, = (8).

This equation, as it stands, gives the changes in period in terms of the
changes of density supposed to be known. And it shews conversely that a
variation of density may always be found which will give prescribed arbitrary
displacements to all the periods. This is a point of some interest.

In order to secure a reasonable continuity in the density, it is necessary to
suppose that \(Sj, \ldots \) are so prescribed that a, assumes ultimately a
constant value when r is increased indefinitely. If this condition be
satisfied, we may take Ao = and then Ar tends to zero as r increases.

As a simple example, suppose that it be required so to vary the density of a
string that, while the pitch of the fundamental tone is displaced, all other
tones shall remain unaltered. The conditions give a,sa, = a =saap = 0.
Accordingly Aq = A = As =0, and .   Thus by (6)
Sp po = - 2ai cos (2wx l).]

141. The differential equation determining the motion of a string, whose
longitudinal density p is variable, is dy d y

P ao =T7, +5 1 dat (1),

from which, if we assume y c (nf), we obtain to determine the normal functions

oY + toy = 0(2),

where v is written for n'+ 7. This equation is of the second order and linear,
but has not hitherto been solved in finite terms. Considered as defining the
curve assumed by the string in the normal mode under consideration, it
determines the curvature at any point, and accordingly embodies a rule by which
the curve can be constructed graphically. Thus in the application to a string
fixed at both ends, if we start from either end at an arbitrary inclination,
and with zero curvature, we are always directed by the equation with what
curvature to proceed, and in this way we may trace out the entire curve.

If the assumed value of v be right, the curve will cross the axis of at the
required distance, and the law of vibration will be completely determined. If
be not known, different values may be tried until the curve ends rightly; a
sufficient approximation to the value of may usually be arrived at by a
calculation founded on an assumed type (Article 88, 90).

Whether the longitudinal density be uniform or not, the periodic time of any
simple vibration varies ceteris paribus as the square root of the density and
inversely as the square root of the tension under which the motion takes place.

The converse problem of determining the density, when the period and the
type of vibration are given, is always soluble. For this purpose it is
only necessary to substitute the given value of y, and of its second
differential coefficient in equation (2). Unless the density be
infinite, the extremities of a string are points of zero curvature,

When a given string is shortened, every component tone is raised in
pitch. For the new state of things may be regarded as derived from the
old by introduction, at the proposed point of fixture, of a spring
(without inertia), whose stiffness is gradually increased without limit.
At each step of the process the potential energy of a given deformation
is augmented, and therefore (Article 88) the pitch of every tone is
raised. In like manner an addition to the length of a string depresses
the pitch, even though the added part be destitute of inertia,

142. Although a general integration of equation (2) of Article 141 is
beyond our powers, we may apply to the problem some of the many
interesting properties of the solution of the linear equation of the
second order, which have been demonstrated by MM. Sturm and Liouville'.
It is impossible in this work to give anything like a complete account
of their investigations; but a sketch, in which the leading features are
included, may be found interesting, and will throw light on some points
connected with the general theory of the vibrations of continuous
bodies. I have not thought it necessary to adhere very closely to the
methods adopted in the original memoirs.

At no point of the curve satisfying the equation
can both y and wy vanish together. By successive differentiations
of (1) it is easy to Prove that, if y and ay vanish simultaneously,
all the higher differential coefficients etc. must also
vanish at the same point, and therefore by Taylor's theorem the curve
must coincide with the axis of x.

Whatever value be ascribed to v , the curve satisfying (1) is sinuous;
being concave throughout towards the axis of 2, since
p is everywhere positive, If at the origin y vanish, and w be positive,
the ordinate will remain positive for all values of x below a certain
limit dependent on the value ascribed to v . - If be very small, the
curvature is slight, and the curve will remain on the positive side of
the axis for a great distance. We have now to prove that as increases, all the
values of x which satisfy the equation y=0 gradually diminish in magnitude.

Let y' be the ordinate of a second curve satisfying the equation

ag well as the condition that y' vanishes at the origin, and let us
suppose that v'' is somewhat greater than v . Multiplying (2) by y,

Dangling\footnote{ The memoirs referred to in the text are contained in the first volume
of Liouville's Journal (1836),}

and (1) by y', subtracting, and integrating with respect to x between
the limits 0 and 2, we obtain, since y and y' both vanish with a, Z; ,
di , i yf -yDe Pr) (3). If we further suppose that 2 corresponds to a
point at which y vanishes, and that the difference between vy'' and is
very small, we get ultimately

id. Y BaP py deren A)

The right-hand member of (4) being essentially positive,
we learn that y' and y are of the same sign, and therefore that,
whether dy be positive or negative, y' is already of the same sign
as that to which y is changing, or in other words, the value of x for
which y' vanishes is less than that for which y vanishes.

If we fix our attention on the portion of the curve lying between = 0
and =1, the ordinate continues positive throughout as the value of
increases, until a certain value is attained, which we will call ,. The
function y is now identical in form with the first normal function u, of
a string of density p fixed at 0 and J, and has no root except at those
points. As again increases, the first root moves inwards from x=/ until,
when a second special value y,' is attained, the curve again crosses the
axis at the point z=/, and then represents the second normal function
u,. This function has thus one internal root, and one only. In like
manner corresponding to a higher value ,' y we obtain the third normal
function u, with two internal roots, and soon. The n function wu, has
thus exactly n -1 internal roots, and since its first differential
coefficient never vanishes simultaneously with the function, it changes
sign each time a root is passed.

From equation (8) it appears that if u, and wu, be two different normal
functions,

t fe 11,16, d= 0 (5).

A beautiful theorem has been discovered by Sturm relating to the number
of the roots of a function derived by addition from a finite number of
normal functions. If u,, be the component of lowest order, and wu, the
component of highest order, the function

f(x) =(6),
where; ,,:, etc. are arbitrary coefficients, has at least m-1 internal
roots, and at most n-1 internal roots. The extremities at =0 and at =I
correspond of course to roots in all cases. The following demonstration
bears some resemblance to that given by Liouville, but is considerably
simpler, and, I believe, not less rigorous.

If we suppose that f(x) has exactly w internal roots (any number of
which may be equal), the derived function f(x) cannot have less than +1
internal roots, since there must be at least one root of /'' (x) between
each pair of consecutive roots of f(x), and the whole number of roots of
f(x) concerned is w+2. In like manner, we see that there must be at
least yw roots of f(z), besides the extremities, which themselves
necessarily correspond to roots; so that in passing from f(.) to f''()
it is impossible that any roots can be lost. Now

f'' (x) = (7)

as we see by (1); and therefore, since p is always positive, we infer
that (8), has at least y roots.

Again, since (8) is an expression of the same form as f(x), similar
reasoning proves that

has at least internal roots; and the process may be continued to any
extent. In this way we obtain a series of functions, all with internal
roots at least, which differ from the original function f(x) by the
continually increasing relative importance of the components of the
higher orders. When the process has been carried sufficiently far, we
shall arrive at a function, whose form differs as little as we please
from that of the normal function of highest order, viz. u,, and which
has therefore n- 1 internal roots. It follows that, since no roots can
be lost in passing down the series of functions, the number of internal
roots of f(a) cannot exceed n - 1.

The other half of the theorem is proved in a similar manner by
continuing the series of functions backwards from f(z). In this way we
obtain

arriving at last at a function sensibly coincident in form with the
normal function of lowest order, viz. u,, and having therefore m -1
internal roots. Since no roots can be lost in passing up the series from
this function to f(z), it follows that f(.) cannot have fewer internal
roots than m-1; but it must be understood that any number of the m - 1
roots may be equal.

We will now prove that f(x) cannot be identically zero, unless all the
coefficients c vanish. Suppose that c, is not zero. Multiply (6) by pw,,
and integrate with respect to x between the limits 0 and 7. Then by (5)

u,J (x) dx = , Uda (9)

from which, since the integral on the right-hand side is finite, we see
that f(x) cannot vanish for all values of included within the range of
integration.

Liouville has made use of Sturm's theorem to shew how a series of normal
functions may be compounded so as to have an arbitrary sign at all
points lying between 2=0 and w=l. His method is somewhat as follows. .

The values of x for which the function is to change sign being a,b,
c,..., quantities which without loss of generality we may suppose to be
all different, let us consider the series of determinants,

w, (a), u,(b), u, (x) u,(a), u,(6), u,(a), u, (a), u,(b), u(x), etc.

The first is a linear function of u, (x) and u,(x), and by Sturm's
theorem has therefore one internal root at most, which root is evidently
a. Moreover the determinant is not identically zero, since the
coefficient of u,(x), viz. u,(a), does not vanish, whatever be the value
of a. We have thus obtained a function, which changes sign at an
arbitrary point a, and there only internally.

u,(a), u,( ) u,(a), u, (2)


The second determinant vanishes when =a, and when z=), and, since it
cannot have more than two internal roots, it changes sign, when passes
through these values, and there only. The coefficient of u,(2) is the
value assumed by the first determinant when z=, and is therefore finite.
Hence the second determinant is not identically zero.

Similarly the third determinant in the series vanishes and changes sign
when = a, when x=6, and when x =c, and at these internal points only.
The coefficient of u,(x) is finite, being the value of the second
determinant when 2=c.

It is evident that by continuing this process we can form functions
compounded of the normal functions, which shall vanish and change sign
for any arbitrary values of x, and not elsewhere internally; or, in
other words, we can form a function whose sign is arbitrary over the
whole range from z=0 to 2=1.

On this theorem Liouville founds his demonstration of the possibility of
representing an arbitrary function between 2=0 and a=l by a series of
normal functions. If we assume the possibility of the expansion and take

Sle) = 14 (0) + fatty )+ 40 (2) + (10), the necessary values of c,, c,,
etc. are determined by (9), and we find

fla) = u(a) I "ou, (a) f(a) da u,2(2) (11).

If the series on the right be denoted by F(a), it remains to establish
the identity of f(x) and F(z).

If the right-hand member of (11) be multiplied by pu,(x) and integrated
with respect to from =0 to w= 1, we see that

or, as we may also write it,

Fe - f(a) ptt,(2e) dat =0 (12), where u,(2) is any normal function, From
(12) it follows that

\[re - f(a) A,u,(2)+A,r4() +A,r4 (ae) +...p dr= 0\]

(18), where the coefficients A,, A,, etc. are arbitrary.

Now if F(x)-f(x) be not identically zero, it will be possible so to
choose the constants A,, A,, etc. that A, u(.) + A,u,(x)+ has
throughout the same sign as F(x) - f(x), in which case every element of
the integral would be positive, and equation (13) could not be true. It
follows that F()-/(x) cannot differ from zero, or that the series of
normal functions forming the right-hand member of (11) is identical with
f(.) for all values of 2 from = 0 to=1,

The arguments and results of this section are of course applicable to the
particular case of a uniform string for which the normal functions are circular.

[As a particular case of variable density the supposition that is
worthy of notice, § 148 In the notation there adopted m + i=n =p o 3',
(14), and the general solution is y = (15). If the string be fixed at
two points, whose abscissae c, are as r to 1, the frequency equation is
= 1 or
\[\ln r\]
where s denotes an integer. The proper frequencies thus depend only upon
the ratio of the terminal abscissae. By supposing r nearly equal to
unity we may pull back upon the usual formula (§ 24) applicable to a
uniform string. The general form of the normal function is
\[ y = \sin(v) \ln (x X_j) \ln(x J(x) )\]
(17).]

\section{Strings}\label{sect:142a} 142 a.
The points where the string remains at rest, or nodes, are of course
determined by the roots of the normal functions, when the vibrations
are free. In this case the frequency is limited to certain definite
values; but when the vibrations are forced, they may be of any
frequency, and it becomes possible to trace the motion of the nodal
points as the frequency increases continuously.

For example, suppose that the imposed force acts at a single point P of
a string AB, whose density may be variable. So long as the frequency is
less than that of either of the two parts AP, PB (supposed to be held at
rest at both extremities) into which the string is divided, there can be
no (interior) node (Q). Otherwise, that part of the string AQ between
the node Q and one extremity (A), which does not include P, would be
vibrating freely, and more slowly than is possible for the longer
length AP, included between the point P and the same extremity. When the
frequency is raised, so as to coincide with the smaller of those proper
to AP, PB, say AP, a node enters at P and then advances towards A, At
each coincidence of the frequency with one of those proper to the whole
string AB, the vibration identifies itself \(i\)th the corresponding free
vibration, and at each coincidence with 8 frequency proper to AP, or
BP, a new node appears at P, and advances in the first case towards A
and in the second towards B. And throughout the whole sequence of events
all the nodes move outward from P towards A or P.

Thus, if the string be uniform and be bisected at P, there is no node
until the pitch rises to the octave (c of the note (c) of the string. At
this stage two nodes enter at P, and move outwards symmetrically. When
ff' is reached, the mode of vibration is that of the free vibration of
the same pitch, and the nodes are at the two points of trisection. At c
these nodes have moved outwards so far as to bisect AP, PP, and two new
nodes enter at P.
143. When the vibrations of a string are not confined to one plane, it
is usually most convenient to resolve them into two sets executed in
perpendicular planes, which may be treated independently. There is,
however, one case of this description worth a passing notice, in which
the motion is most easily conceived and treated without resolution.

Suppose that
\[y= \sin(r) \cos(Z) z= \sin() sin\]
(1). Then Let (2), and (3), shewing that the whole string is at any
moment in one plane, which revolves uniformly, and that each particle
describes a circle
with radius sin. In fact, the whole system turns without
relative displacement about its position of equilibrium, completing each
revolution in the time \(t+s\). The mechanics of this case is quite as
simple as when the motion is confined to one plane, the resultant of the
tensions acting at the extremities of any small portion of the string's
length being balanced by the centrifugal force.

144. The general differential equation for a uniform string, viz.

Py = ty ae =a dx (1),

may be transformed by a change of variables into (2),
where u=a-at,v=2x+at. The general solution of (2) is
y=f(u) + Fo) =feat) + P(e +at)(8),

SJ, F being two arbitrary functions.
\footnote{[Equations (1) and (3) are due to D'Alembert (1750).]}

Let us consider first the case in which F' vanishes. When t has any
particular value, the equation

y =f (a - at) (4),
expressing the relation between x and y, represents the form of the
string. A change in the value of c is merely equivalent to an alteration
in the origin of x, so that (4) indicates that a certain form is
propagated along the string with uniform velocity a in the positive
direction. Whatever the value of y may be at the point and at the time
c, the same value of y will obtain at the point z + At at the time t+
At.

The form thus perpetuated may be any whatever, so long as it does not
violate the restrictions on which (1) depends.

When the motion consists of the propagation of a wave in the positive
direction, a certain relation subsists between the inclination and the
velocity at any point. Differentiating (4) we find

Initially,

above relation be not satisfied, the motion cannot be represented by
(4). In a similar manner the equation

(6),
denotes the propagation of a wave in the negative direction, and the
relation between dyfdt and dyidx corresponding to (5) is
dt dx (7).
In the general case the motion consists of the simultaneous propagation
of two waves with velocity a, the one in the positive,
and the other in the negative direction; and these waves are
entirely independent of one another. In the first
\( dy/dt = - a dy/dx \)
and in the second
\( dy/dt = a dy/dx \).
The initial values of \( dy/dt\) and \( dy /dx\)  must be
conceived to be divided into two parts, which satisfy respectively the
relations (5) and (7). The first constitutes the wave which will advance
in the positive direction without change of form; the second, the
negative wave. Thus, initially,

whence

equations which determine the functions f' and F'' for all values of the
argument from =- 0 tow= 0, if the initial values of ay and oy be known.

If the disturbance be originally confined to a finite portion of the
string, the positive and negative waves separate after the interval of
time required for each to traverse half the disturbed portion.

Fig. 23.

Q B 0 A P

144. POSITIVE AND NEGATIVE WAVES. 181

Suppose, for example, that AB is the part initially disturbed. A point P
on the positive side remains at rest until the positive wave has
travelled from A to P, is disturbed during the passage of the wave, and
ever after remains at rest. The negative wave never affects P at all.
Similar statements apply, mutatis mutandis, to a point Q on the negative
side of AB. If the character of the
original disturbance be such that ay -3 1 dy vanishes initially, there
is no positive wave, and the point P is never disturbed at all;

dy , l dy dy and if - de ta = p vanish initially, there is no negative
wave. If di
vanish initially, the positive and the negative waves are similar and
equal, and then neither can vanish. In cases where either wave vanishes,
its evanescence may be considered to be due to the mutual destruction of
two component waves, one depending on the initial displacements, and the
other on the initial velocities, On the one side these two waves
conspire, and on the other they destroy one another. This explains the
apparent paradox, that P can fail to be affected sooner or later after
4B has been disturbed.

The subsequent motion of a string that is initially displaced without
velocity, may be readily traced by graphical methods. Since the positive
and the negative waves are equal, it is only necessary to divide the
original disturbance into two equal parts, to displace these, one to the
right, and the other to the left, through a space equal to at, and then
to recompound them. We shall presently apply this method to the case of
a plucked string of finite length.

145. Vibrations are called stationary, when the motion of each particle
of the system is proportional to some function of the time, the same for
all the particles. If we endeavour to satisfy

By ty (L),

4 =a"

at' da by assuming y= XT, where X denotes a function of only, and T a
function of c only, we find 1 T 1 X

Tdi X dx m (a constant), so that

\[X= \cos(m) + D\]

(2) proving that the vibrations must be simple harmonic, though of
arbitrary period. The value of y may be written y =P (mat - e) (mx -a) .
=4P (mat + mz-e -a) +4P (mat -ma-e+a)(3),
shewing that the most general kind of stationary vibration may be
regarded as due to the superposition of equal progressive vibrations,
whose directions of propagation are opposed. Conversely, two stationary
vibrations may combine into a progressive one.

The solution y=/(-at)+F (a+ at) applies in the first instance to an
infinite string, but may be interpreted so as to give the solution of
the problem for a finite string in certain cases. Let us suppose, for
example, that the string terminates at 2=0, and is held fast there,
while it extends to infinity in the positive direction only. Now so long
as the point =0 actually remains at rest, it is a matter of indifference
whether the string be prolonged on the negative side or not. We are thus
led to regard the given string as forming part of one doubly infinite,
and to seek whether and how the initial displacements and velocities on
the negative side can be taken, so that on the whole there shall be no
displacement at 2 = 0 throughout the subsequent motion. The initial
values of y and on the positive side determine the corresponding parts
of the positive and negative waves, into which we know that the whole
motion can be resolved. The former has no influence at the point =0. On
the negative side the positive and the negative waves are initially at
our disposal, but with the latter we are not concerned. The problem is
to determine the positive wave on the negative side, so that in
conjunction with the given negative wave on the positive side of the
origin, it shall leave that point undisturbed.

Fig. 24.

Let OP QRS... be the line (of any form) representing the wave in OX,
which advances in the negative direction. It is
evident that the requirements of the case are met by taking on the other
side of O what may be called the contrary wave, so that O is the
geometrical centre, bisecting every chord (such as PP') which passes
through it. Analytically, if y=/(zx) is the equation of y=f(.) is the
equation of

When after a time c the curves are shifted to the left and to the right
respectively through a distance at, the coordinates corresponding to
x=0Q are necessarily equal and opposite, and therefore when compounded
give zero resultant displacement.

The effect of the constraint at O may therefore be represented by
supposing that the negative wave moves through undisturbed, but that a
positive wave at the same time emerges from O. This reflected wave may
at any time be found from its parent by the following rule :

Fig. 25.

Let APQRS... be the position of the parent wave. Then the reflected wave
is the position which this would assume, if it were
turned through two right angles, first about OX as an axis of rotation,
and then through the same angle about OY. In other words, the return
wave is the image of APQRS formed by successive optical reflection in OX
and OY, regarded as plane mirrors. ;

The same result may also be obtained by a more analytical process. In
the general solution

y =f (- at) + F(x + at), the functions f(z), F (z) are determined by the
initial circumstances for all positive values of z. The condition at 2=0
requires that

J (- at) + (F (at) =0 for all positive values of c, or

f(-2)=- Fe)

for positive values of z The functions f and F are thus determined for
all positive values of and c.

There is now no difficulty in tracing the course of events when two
points of the string A and B are held fast. The initial disturbance in
AB divides itself into positive and negative waves, which are reflected
backwards and forwards between the fixed points, changing their
character from positive to negative, and vice versa, at each reflection.
After an even number of reflections in each case the original form and
motion -is completely recovered. The process is most easily followed in
imagination when the initial disturbance is confined to a small part of
the string, more particularly when its character is such as to give rise
to a wave propagated in one direction only. The pulse travels with
uniform velocity (a) to and fro along the length of the string, and
after it has returned a second time to its starting point the original
condition of things is exactly restored. The period of the motion is
thus the time required for the pulse to traverse the length of the
string twice, or

The same law evidently holds good whatever may be the character of the
original disturbance, only in the general case it may happen that the
shortest period of recurrence is some aliquot part of 7.

146. The method of the last few sections may be advantageously applied
to the case of a plucked string. Since the initial velocity vanishes,
half of the displacement belongs to the positive and half to the
negative wave. The manner in which the wave must be completed so as to
produce the same effect as the constraint, is shewn in the figure, where
the upper curve represents

Fig. 26, SS P Pp P 3 As As 4 Ae Ag Q Q Q P Pp P As 7 Ag A4 Ag Ag Q Q Q a

the positive, and the lower the negative wave in their initial
positions. In order to find the configuration of the string at any

future time, the two curves must be superposed, after the upper has been
shifted to the right and the lower to the left through a space equal to
at.

146. - GRAPHICAL METHOD. 185

The resultant curve, like its components, is made up of straight pieces.
A succession of six at intervals of a twelfth of the period,

Fig. 27.

shewing the course of the vibration, is given in the figure (Fig. 27),
taken from Helmholtz. From G the string goes back again to A through the
same stages
\footnote{This method of treating the vibration of a plucked
string is due to Young. Phil. Trans., 1800. The student is recommended
to make himself familiar with it by actually constructing the forms of Fig. 27.}.

It will be observed that the inclination of the string at the points of
support alternates between two constant values.

147. If a small disturbance be made at the time c at the point x of an
infinite stretched string, the effect will not be felt at O until after
the lapse of the time +a, and will be in all respects the same as if a
like disturbance had been made at the point x + x at time t - x/a.
Suppose that similar disturbances are communicated to the string at
intervals of time r at points whose distances from O increase each time
by a \(\delta \tau\), then it is evident that the result at O will be
the same as if the disturbances were all made at the same point,
provided that the time intervals be increased from 7 to r +87. This
remark contains the theory of the alteration of pitch due to motion of
the source of disturbance; a subject which will come under our notice
again in connection with aerial vibrations.

148. When one point of an infinite string is subject to a forced
vibration, trains of waves proceed from it in both directions according
to laws, which are readily investigated. We shall suppose that the
origin is the point of excitation, the string being there subject to the
forced motion y= Ae''; and it will be sufficient to consider the
positive side. If the motion of each element ds be resisted by the
frictional force yds, the differential equation is

d. dy ,d Gate a= atk (1); or since y e, d CK: T= ry (2),

if for brevity we write 0 for the coefficient of y. The general solution
is (8).

Now since y is supposed to vanish at an infinite distance, D must
vanish, if the real part of X be taken positive. Let

where is positive. Then the solution is y = Aen (ar B) a time (4), or,
on throwing away the imaginary part, y= Ae (pt-BL) (5),
corresponding to the forced motion at the origin Y= A (Pb... (6)

An arbitrary constant may, of course, be added to c.

To determine a and 8, we have

a-p=


If we suppose that is small, =+, a=

a 2a and y = Ae 4'' (pt-f) (8).

This solution shews that there is propagated along the string a wave,
whose amplitude slowly diminishes on account of the exponential factor.
If =0, this factor disappears, and we have simply

y = A (ot - Pe) se settee (9).

This result stands in contradiction to the general law that, when there
is no friction, the forced vibrations of a system
(due to a single simple harmonic force) must be synchronous in phase throughout.
According to (9), on the contrary, the phase varies continuously in
passing from one point to another along the string. The fact is, that we
are not at liberty to suppose =0 in (8), inasmuch as that equation was
obtained on the assumption that the real part of in (3) is positive, and
not zero. However long a finite string may be, the coefficient of
friction may be taken so small that the vibrations are not damped before
reaching the further end. After this point of smallness, reflected waves
begin to complicate the result, and when the friction is diminished
indefinitely, an infinite series of such must be taken into account, and
would give a resultant motion of the same phase throughout.

This problem may be solved for a string whose mass is supposed to be
concentrated at equidistant points, by the method of Article 120. The
coordinate y, may be supposed to be given (= Ae), and it will be found
that the system of equations (5) of Article 120 may all be satisfied by
taking

where 6 is a complex constant determined by a quadratic equation. The
result for a continuous string may be afterwards deduced.

{[}In the notation of § 120 the quadratic equation is + = 0 where
A = fp +
The roots of (11) are (11).

(12). 2B (13) and are imaginary if A , that is, if fia .(14),
a condition always satisfied in passing to the limit where a and fi are
infinitely small. In any case when (14) is satisfied the modulus of 0 is
unity, so that (10) represents wave propagation.

If, however, (14) be not satisfied, the values of 0 are real. In this
case all the motions are in the same phase, and no wave is propagated.
The vibration impressed upon is imitated upon
a reduced scale by with amplitudes which form a
geometrical progression. In the first case the motion is propagated to
an infinite distance, but in the second it is practically confined to a
limited region round the source.{]}

\section{148 a}\label{sect:148a}
148 a.
So long as the conditions of § 144 are satisfied, a positive, or a
negative, wave is propagated undisturbed. If however there be any want
of uniformity, such (for example) as that caused by a load attached at a
particular point, reflection will ensue when that point is reached. The
most interesting problem under this head is that of two strings of
different longitudinal densities, attached to one another, and vibrating
transversely under the common tension.  Or, if we regard the string as
single, the density may be supposed to vary discontinuously from one
uniform value (p ) to another (pj) If Oi, Os denote the corresponding
velocities of propagation. ai =ri pi, = (1),
and i = Oi o, = a ( Pi) (2). The conditions to be satisfied at the
junction of the two parts are (i) the continuity of the displacement y,
and (ii) the continuity of dy dx. If the two parts met at a finite
angle, an infinitely small element at the junction would be subject to a
finite force.

Let us suppose that a positive wave of harmonic type, travelling in the
first part (pO, impinges upon the second (p ). In the latter the motion
will be adequately represented by a positive wave, but in the former we
must provide for a negative 'reflected wave. Thus we may take for the
two parts respectively y = (3),

y=Z w (4),

where = Sw Xi, kn = 2ir ,

so that kiOi = k (5) The conditions at the junction (x= 0) give

H+K=L whence kiH  kiK = hjj

K kih t1 H ib, M + 1 (6),

0)

.(8). Since the ratio KjH is real, we may suppose that both quantities
are real; and if we throw away the imaginary parts from (3) and (4) we
get as the solution in terms of real quantities y H cos ki (Oit x) + K
cos (oyt + a) (9) ;

y(H + K) cos k (a4x) (10). The ratio of amplitudes of the reflected and
the incident waves expressed by (8) is that first obtained b) T. Young
for the corresponding problem in Optica

\section{148 b}\label{sect:148b}

The expression for the intensity of reflection established in § 148 a
depends upon the assumption that the transition from the one density to
the other is sudden, that is occupies a distance which is small in
comparison with a wave length. If the transition be gradual, the
reflection may be expected to fall off, and in the limit to disappear
altogether.

The problem of gradual transition includes, of course, that of a
variable medium, and would in general be encumbered with great
difficulties There is, however, one case for which the solution may be
readily expressed, and this it is proposed to consider in the present
section. The longitudinal density is supposed to vary as the inverse
square of the abscissa. If y, denoting the transverse displacement be
proportional to the equation which it must satisfy as a function of a,
is (§ 141), + nar y 0 (1),

where n is some positive constant, of the nature of an abstract number.

The solution of (1) is y = (2),

where i (3X

If m be real, that is, if \(n > n\) , we may obtain, by supposing 0, as
a final solution in real quantities,

y as Ca ooe (pt m log a + e) (4),

which represents a positive progressive wave, in many respects similar
to those propagated in uniform media.

Let us now suppose that, to the left of the point the
valuable medium is replaced by one of uniform constitution, such that
there is no discontinuity of density at the point of transition ; and
let us inquire what reflection a positive progressive wave in the
uniform medium will undergo on arrival at the variable medium. It will
be sufficient to consider the case where m is real, that is, where the
change of density is but moderately rapid.

By supposition, there is no negative wave in the variable medium, so
that il = 0 inX2). Thus

y = = (i ;

and, when 07 = = (5).

y dx x ' '

The general solution for the uniform medium, satisfying the equation + n
xr = 0, may be written from which, when x + Ke (6)i dy in HK yd ' Xj
.(7). In equation (6), H represents the amplitude of the incident
positive wave, and K the amplitude of the reflected negative wave. The
condition to be satisfied at = is expressed by
equating the values of given by (5) and (7). Thus
which gives, in symbolical form, the ratio of the reflected to the
incident vibration having regard to (3), we may write (8) in the form

so that the amplitude of the reflected wave is (n + w) of that of the
incident. Thus, as was to be expected, when n and m are great, i.e.,
when the density changes slowly in the variable
medium, there is but little reflection. As regards phase, the result
embodied in (9) may be represented by supposing that the reflection
occurs = and involves a change of phase amounting to a quarter period.

Passing on now to the more important problem, we will suppose that the
variable medium extends only so far as the point beyond which the
density retains uniformly its value at that point. A positive wave
travelling at first in a uniform medium of density proportional to
passes at the point a = ai into a variable medium of density
proportional to or , and again, at the point x , into a uniform medium
of density proportional to The velocities of propagation are inversely
proportional to the square roots of the densities, so that, if be the
refractive index between the extreme media, The thickness (d) of the
layer of transition is

d  x Xi

The wave-lengths in the two media are given by

that

For the first medium we take, as before. if = He . a5i  +in

+ Ke giving, when dy mH K y  XiH+K Xi TT jr

for brevity, we write 6 tat .  Jti K

For the variable medium. +

giving, when x = a

dy + tm) Ax + (j ) Bx (10).  (11). (12): (6).  (7). ..(2), (13).

These expressions allow ns to trace the effect of a more or less gradual
reduction between media of given indices. If the transition be
absolutely abrupt, n=sO, by (12); so that In this case, (18) gives us
(§ 148 a) Young's well-known formula it (21). sin li cc

Since increases continually from = the ratio (19)

increases continually from m' = 0 to m' = J, is., diminishes continually
from the case of sudden transition m' = when its value is (21), to the
critical case m = 0, when its value is (20), after which this form no
longer holds good. When rn! = 0, n = and, by (12), d (X,  )l 4sw. When
is the appropriate form. We see from it that with increasing n the
reflection diminishes, until it vanishes, when mlogfisTT, i.6. when

(log a)(22). With a still more gradual transition the reflection
revives, reaches a maximum, again vanishes when m log ft == 2ir, and so
on \footnote{Proc. Math. Soc., vol. XI, February, 1880; where will also be found a
numerical example illustrative of optical conditions.}

\section{148 c}\label{sect:148c}

In the problem of connected strings, vibrating under the influence of
tension alone, the velocity in each uniform part is independent of wave
length, and there is nothing corresponding to optical dispersion. This
state of things will be departed from if we introduce the consideration
of stiffness,and it may be of interest to examine in a simple case how
far the problem of reflection is thereby modified. As in §148 a, we will
suppose that at d; = 0 the density changes discontinuously from pi to
p2i but that now the vibrations of the second pieut occur under the
influence of sensible stiffness. The differential equation applicable in
this case is, § 188,

`` d da diP cr, if y vaiy as that, if y vary as e . (1). 8 lo, n = 0, (2). '

In consequence of the stiffness represented by the velocity of
propagation deviates from Os, and must be found from (2). The two values
of given by this equation are real, one being positive and the other
negative. The four admissible values of k may thus
be written + A:,, ± so that the complete solution of (1) will be

y = + Cer + (3),

Aa, fca being real and positive. The velocity of propagation is

In the application which we have to make the disturbance of the
imperfectly flexible second part is due to a positive wave entering it
from the first part. When x is great and positive, (3) must reduce to
its second term. Thus D = 0;

and we are left with y = Be + Ce (4).

This holds when x is positive. When x is negative, corresponding to the
perfectly flexible first part, we have y 5= ....
(5).

in which = n oj
(6)

(7).

The conditions at the junction are first the continuity of y and dyfdx.
Further, d yfda in (4) must vanish at this place, inasmuch as curvature
implies a couple (§ 162), and this could not bu transmitted by the first
part. Hence

H+K B+G (8).

k,(HK) = k ih,C ..  (9).

Jbi 5 + A, O=0 . (10).

From these we deduce

H + K k,(h, + k,) HK kji (11).

K(h)+ii H A + (12) :

and thence for the intended of reflection, equal to Mod'. (KjH),

( , + ,) +

If the second part, as well as the first, be perfectly flexible, = 0, Ai
= 00 , and we

back on Young's formula. In general, the intensity of reflection is not
accurately given by this formula, even though we employ therein the
value of the refractive index appropriate to the waves actually under
propagation.
